\documentclass[12pt,a4paper]{article}

\usepackage[utf8]{inputenc}
\usepackage[T1]{fontenc}
\usepackage{amsmath,amssymb,amsfonts}
\usepackage{mathtools}
\usepackage{booktabs}
\usepackage{array}
\usepackage{graphicx}
\usepackage{hyperref}
\usepackage[margin=1in]{geometry}
\usepackage{caption}
\usepackage{float}
\usepackage{enumitem}
\usepackage{xcolor}
\usepackage{cite}

\hypersetup{
  colorlinks=true,
  linkcolor=blue!70!black,
  citecolor=blue!70!black,
  urlcolor=blue!70!black
}

\newcommand{\vp}{\varphi}
\newcommand{\aGUT}{\alpha_{\mathrm{GUT}}}
\newcommand{\MPl}{M_{\mathrm{Pl}}}
\newcommand{\mZ}{m_Z}
\newcommand{\QGUT}{Q_{\mathrm{GUT}}}
\newcommand{\SM}{\mathrm{SM}}
\newcommand{\MSSM}{\mathrm{MSSM}}

\title{A Golden-Ratio Lattice for Dimensionless Coupling Constants:\\
Structure, Predictions, and Falsification}

\author{Tristan~Stevens\\[4pt]
\textit{Independent Researcher}\\
\texttt{[email placeholder]}}

\date{\today}

\begin{document}
\maketitle

% ===================================================================
%  ABSTRACT
% ===================================================================
\begin{abstract}
We investigate the hypothesis that dimensionless parameters of the Standard
Model organize on a discrete lattice $G(C,m) = C/\vp^m$, where
$\vp = (1+\sqrt{5})/2$ is the golden ratio, $m$ is an integer, and $C$
belongs to a short menu $\{15, 45, 60, 120, 180, 360\}$ derived from
Lie-algebra invariants of the Standard Model gauge groups
$U(1)\times SU(2)\times SU(3)$.  All model rules---integer stepping,
inverse-coupling orientation, gauge-derived $C$ menu, target definitions, and
tolerance thresholds---are frozen before any data comparison.  Across 43
independent dimensionless ratios spanning gauge couplings, fermion masses,
mixing angles, anomalous magnetic moments, and cosmological parameters, we
observe a sub-percent hit rate of 42\%, corresponding to $2.68\times$
enrichment over a null model ($p < 5\times 10^{-5}$).  Lattice-constrained
gauge coupling unification in the MSSM converges on a single lattice point
$(C=15,\,m=-1)$ with $\aGUT^{-1} = 15\vp \approx 24.27$, where $C=15 =
360/\dim(\mathrm{SU}(5))$.  The Planck-to-proton mass hierarchy and the
cosmological constant discrepancy reduce to single lattice integers:
$\MPl/m_p = \vp^{91}$ ($+0.51\%$) and $\rho_{\mathrm{Pl}}/\rho_{\mathrm{vac}}
= \vp^{588}$ ($+0.01\%$).  We present 10 registered, falsifiable predictions:
particle-physics observables testable by DUNE, JUNO, Hyper-Kamiokande, and
FCC-ee---including $\delta_{CP}^{\mathrm{PMNS}} = 201.2^\circ$,
$\sin^2\theta_{13} = 0.02077$ (currently in $2.3\sigma$ tension), and a proton
lifetime of $\sim\!10^{34\text{--}35}$~years---together with three vacuum
spectral predictions (dynamic Casimir photon peaks at $\vp$-harmonic
frequencies, anomalous Casimir force at lattice plate separations, and excess
noise in a cryogenic 560~GHz cavity) that directly probe the cosmological
constant problem.  The companion open-source codebase reproduces every
numerical result reported here.
\end{abstract}

\tableofcontents
\newpage

% ===================================================================
%  1. INTRODUCTION
% ===================================================================
\section{Introduction}
\label{sec:intro}

The Standard Model of particle physics contains roughly two dozen free
parameters---coupling constants, fermion masses, mixing angles, and a CP
phase---whose values are measured but not explained.  Why the fine-structure
constant $\alpha \approx 1/137$?  Why the proton-to-electron mass ratio
$m_p/m_e \approx 1836$?  These questions have motivated theoretical programs
ranging from grand unification~\cite{georgi1974,pati1974} and string
landscape statistics~\cite{bousso2000} to anthropic reasoning~\cite{barrow1986}.
A complementary, more modest question is: \emph{do the dimensionless parameters
exhibit a shared discrete structure that could point toward an underlying
organizing principle?}

Attempts to find numerical patterns among fundamental constants have a long
and largely cautionary history.  Eddington's algebraic derivations of
$1/\alpha = 136$ (later revised to~137) and Dirac's large-number
hypothesis~\cite{dirac1937} are classic examples of numerological programs
that, while intellectually stimulating, lacked falsifiability and ultimately
did not survive experimental scrutiny.  The lesson from these efforts is not
that pattern-seeking is illegitimate, but that it must be accompanied by
\emph{frozen rules}, \emph{pre-registered predictions}, and \emph{honest
failure modes}.

In this paper, we describe a minimal framework in which dimensionless Standard
Model parameters are compared against a discrete lattice
\begin{equation}
  G(C,m) = \frac{C}{\vp^m}, \qquad m \in \mathbb{Z},
  \label{eq:lattice}
\end{equation}
where $\vp = (1+\sqrt{5})/2 \approx 1.618$ is the golden ratio and $C$ is
restricted to six values derived from the Lie-algebra invariants of
$U(1)\times SU(2)\times SU(3)$.  The framework is not a theory---it provides
no dynamical mechanism, Lagrangian, or derivation from first principles.  It
is a \emph{phenomenological hypothesis under stress-test}: a claim that an
unexpectedly large fraction of measured dimensionless ratios land near points
on this lattice, at a level that appears to exceed chance.

We emphasize what this work is \emph{not}:
\begin{itemize}[nosep]
  \item It is not a replacement for the Standard Model or any of its
    precision calculations.
  \item It is not a derivation of coupling constants from first principles.
  \item It does not explain \emph{why} the golden ratio or the base~360
    should play any role.
\end{itemize}
What it \emph{is}: a transparent, reproducible audit of a specific discrete
hypothesis, with all rules frozen before data comparison, yielding falsifiable
predictions testable by current and near-future experiments.

The paper is organized as follows.  Section~\ref{sec:model} defines the model
and its anti-overfitting contract.  Section~\ref{sec:results} presents results
across gauge couplings, mass ratios, mixing parameters, and gravity.
Section~\ref{sec:stats} assesses statistical significance.
Section~\ref{sec:finetuning} reframes the hierarchy and cosmological constant
problems as lattice integers.  Section~\ref{sec:predictions} presents
registered falsifiable predictions.  Section~\ref{sec:limitations} discusses
limitations and open questions.  Section~\ref{sec:conclusion} concludes.


% ===================================================================
%  2. THE MODEL
% ===================================================================
\section{The Model}
\label{sec:model}

\subsection{Core definitions}
\label{sec:core}

The golden ratio is defined as
\begin{equation}
  \vp \equiv \frac{1+\sqrt{5}}{2} \approx 1.6180339887\,.
\end{equation}
The phi-lattice assigns to each pair $(C, m)$ the value
\begin{equation}
  G(C,m) = \frac{C}{\vp^m}, \qquad m \in \mathbb{Z}\,.
  \label{eq:lattice_def}
\end{equation}
A physical dimensionless quantity $X$ is said to ``hit'' the lattice at
tolerance $\epsilon$ if there exists a pair $(C,m)$ in the allowed menu such
that
\begin{equation}
  \left|\frac{G(C,m) - X}{X}\right| \le \epsilon\,.
\end{equation}

\subsection{The $C$-menu from gauge-group invariants}
\label{sec:cmenu}

The single free structural choice in the model is the base number.  We adopt
$\mathrm{base} = 360$ and generate $C$ candidates by dividing by
representation-independent Lie-algebra invariants of the Standard Model gauge
factors: the dimension $\dim(\mathfrak{g})$, the Coxeter number~$h$, the dual
Coxeter number~$h^\vee$, and their pairwise products.

\begin{table}[h]
\centering
\caption{Lie-algebra invariants of the SM gauge groups.}
\label{tab:invariants}
\begin{tabular}{@{}lccc@{}}
\toprule
Group & Rank~$r$ & $\dim(\mathfrak{g})$ & $h = h^\vee$ \\
\midrule
$U(1)$  & 1 & 1 & --- \\
$SU(2)$ & 1 & 3 & 2   \\
$SU(3)$ & 2 & 8 & 3   \\
\bottomrule
\end{tabular}
\end{table}

The construction rules are:
\begin{align}
  C &= \mathrm{base}, \quad
  C = \frac{\mathrm{base}}{\dim}, \quad
  C = \frac{\mathrm{base}}{h}, \quad
  C = \frac{\mathrm{base}}{h^\vee}, \quad
  C = \frac{\mathrm{base}}{\dim \cdot h}\,.
\end{align}
Applied to each gauge factor and de-duplicated, this yields the menu
\begin{equation}
  \boxed{C \in \{15,\; 45,\; 60,\; 120,\; 180,\; 360\}}\,.
  \label{eq:cmenu}
\end{equation}
Explicitly: $U(1)$ contributes $\{360\}$; $SU(2)$ contributes
$\{120, 180, 60\}$; $SU(3)$ contributes $\{45, 120, 15\}$.  After
de-duplication, six values remain.

The base~360 is selected as the smallest highly-composite integer whose
gauge-derived $C$~set is short (six values after de-duplication) and contains
hits for the EM, strong, and weak anchor targets at integer~$m$ within 5\%
tolerance.  Its mathematical properties---$360 = 2^3 \times 3^2 \times 5$
with 24~divisors, and the conventional representation of a full
$2\pi$~cycle---make it a natural starting point.

\subsection{Anti-overfitting contract}
\label{sec:contract}

The following rules are frozen before any data comparison:

\begin{enumerate}[nosep, label=\textbf{F\arabic*.}]
  \item \textbf{Integer stepping:} $m \in \mathbb{Z}$ everywhere; no
    real-valued $m$ fitting.
  \item \textbf{Inverse-coupling orientation:} all targets use
    $1/\alpha$-style (inverse coupling) convention.
  \item \textbf{Gauge-derived $C$~menu:} only Eq.~\eqref{eq:cmenu};
    no ad~hoc additions.
  \item \textbf{Fixed targets:} specific dimensionless quantities defined
    before testing (Section~\ref{sec:targets}).
  \item \textbf{Fixed tolerance:} 5\% for strict anchor tests; 2\% for
    out-of-sample RG cross-checks.
  \item \textbf{No post-hoc rescue:} if a target fails, it stays failed.
    Changing the $C$~menu or target definition to rescue a failure requires
    re-running all pre-registered tests under the modified contract.
\end{enumerate}

\subsection{Target definitions}
\label{sec:targets}

\textbf{Electromagnetic:} the low-energy fine-structure inverse
$1/\alpha = 137.036$.

\textbf{Strong:} $1/\alpha_s(m_H)$, obtained by 1-loop QCD running from
$\alpha_s(\mZ) = 0.1179$ to $m_H = 125$~GeV with $n_f = 5$, yielding
$1/\alpha_s(m_H) \approx 8.87$.

\textbf{Weak:} $1/\alpha_2$, where $\alpha_2 = \alpha(\mZ)/\sin^2\theta_W^{\mathrm{OS}}$
and $\sin^2\theta_W^{\mathrm{OS}} = 1 - m_W^2/m_Z^2$.  This gives
$1/\alpha_2 \approx 28.54$.

\textbf{Hypercharge (GUT-normalized):}
$1/\alpha_1^{\mathrm{GUT}} = (3/5)\,\alpha_Y^{-1}$
with $\alpha_Y = \alpha(\mZ)/\cos^2\theta_W^{\mathrm{OS}}$, yielding
$1/\alpha_1^{\mathrm{GUT}} \approx 59.65$.

\textbf{Gravity:} $1/\alpha_G(M) = \hbar c/(G_N M^2)$ with frozen mass
anchors: proton ($m_p$), electron ($m_e$, cross-check), and
gravitational-wave band anchors (LIGO, LISA, PTA, CMB).

All external inputs---$\alpha(\mZ)$, $\alpha_s(\mZ)$, $m_W$, $m_Z$,
$m_H$, $G_N$, particle masses---are taken from the Particle Data
Group~\cite{pdg2024}.


% ===================================================================
%  3. RESULTS
% ===================================================================
\section{Results}
\label{sec:results}

\subsection{Core anchor hits}
\label{sec:anchors}

Under the frozen $C$~menu and integer~$m$, the four gauge-coupling inverse
targets yield:

\begin{table}[H]
\centering
\caption{Strict lattice fits to SM gauge couplings.}
\label{tab:anchors}
\begin{tabular}{@{}llrrrr@{}}
\toprule
Sector & Target & Value & $(C, m)$ & $C/\vp^m$ & Rel.~err \\
\midrule
EM & $1/\alpha$ & 137.036 & $(360, 2)$ & 137.508 & $+0.34\%$ \\
Strong & $1/\alpha_s(m_H)$ & 8.867 & $(60, 4)$ & 8.754 & $-1.27\%$ \\
Weak & $1/\alpha_2$ & 28.536 & $(120, 3)$ & 28.328 & $-0.73\%$ \\
Hypercharge & $1/\alpha_1^{\mathrm{GUT}}$ & 59.652 & $(60, 0)$ & 60.000 & $+0.58\%$ \\
\bottomrule
\end{tabular}
\end{table}

All four inverse couplings are captured within 1.3\% by lattice points using
only three distinct $C$~values (60, 120, 360) and $m$~values in the range
$[0, 4]$.  The EM hit $360/\vp^2 \approx 137.5$ is the motivating observation;
the strong, weak, and hypercharge hits are independent confirmations.

The force hierarchy is encoded in the $m$-ordering: strong ($m=4$), weak
($m=3$), EM ($m=2$), hypercharge ($m=0$).  Adjacent forces differ by
$\Delta m = 1$--$2$ steps, each multiplying the coupling by
$\vp^{\Delta m} \approx 1.6$--$2.6$.

\subsection{RG-within-band validation}
\label{sec:rg}

The lattice assigns each coupling a single anchor $(C,m)$.  Standard
renormalization group (RG) running then propagates each coupling to other
energy scales \emph{with no additional fitting}.  This tests whether the
lattice is consistent with the known scale dependence.

For each gauge sector, we fit the anchor once, then RG-evolve to $m_W$,
$m_H$, $m_t$, 1~TeV, and 10~TeV.  Results:

\begin{itemize}[nosep]
  \item \textbf{Strong:} $1/\alpha_s(Q)$ passes at 2\% across all
    cross-check scales (typical errors 1--2\%).
  \item \textbf{EM:} $1/\alpha \to 1/\alpha(\mZ)$ passes at 2\% under
    deterministic QED running.
  \item \textbf{Weak:} $1/\alpha_2(Q)$ passes at 2\% across $m_W$, $m_H$,
    1~TeV, 10~TeV.
  \item \textbf{Hypercharge:} $1/\alpha_1^{\mathrm{GUT}}(Q)$ passes at 2\%
    across the same scales.
  \item \textbf{Electroweak mixing:} the derived
    $\sin^2\theta_W(Q) = \alpha_2^{-1}/(\alpha_2^{-1} + \alpha_1^{-1})$
    passes at 2\% with typical offset~$\sim\!1\%$.
\end{itemize}

The interpretation is that $m$ serves as a \emph{coarse band index},
selecting the approximate coupling strength, while standard RG flow provides
the continuous within-band motion.

\subsection{Gravity on the lattice}
\label{sec:gravity}

The gravitational coupling $\alpha_G(M) = G_N M^2/(\hbar c)$ is
mass-dependent, motivating distinct ``gravity types'' for different mass
anchors.  Under the frozen $C$~menu:

\begin{table}[H]
\centering
\caption{Gravity fits across mass anchors.}
\label{tab:gravity}
\begin{tabular}{@{}llrrr@{}}
\toprule
Gravity type & Mass anchor & $(C, m)$ & Rel.~err & Character \\
\midrule
Ordinary (proton)  & $m_p$                   & $(45, -175)$  & $-0.61\%$ & Macro \\
Cross-check (electron) & $m_e$               & $(360, -202)$ & $+3.6\%$  & Macro \\
LIGO band          & $5.4\times 10^{13}$~GeV & $(360, -39)$  & $<0.01\%$ & Stellar \\
LISA band          & $2.2\times 10^{12}$~GeV & $(15, -59)$   & $<0.01\%$ & SMBH \\
PTA band           & $1.6\times 10^{9}$~GeV  & $(15, -89)$   & $<0.01\%$ & Cosmic \\
CMB primordial     & $2.9\times 10^{4}$~GeV  & $(45, -132)$  & $<0.01\%$ & Inflationary \\
Planck             & $m_P$                   & $(120, 10)$   & $-2.4\%$  & Quantum \\
\bottomrule
\end{tabular}
\end{table}

Planck-scale gravity sits at positive~$m$ ($m \approx +10$), within the gauge
cluster ($m = 0$ to $+4$).  Ordinary-matter gravity sits at large
negative~$m$ ($m = -175$).  The negative $m$-axis forms a gradient from
local gravitational phenomena ($m \sim -39$) to cosmic scales ($m \sim -132$).
The hierarchy problem---why gravity appears $\sim\!10^{40}$ times weaker than
EM---becomes a count of $\sim\!180$ lattice steps separating their respective
$m$-addresses.

The electron--proton cross-check is self-consistent: the predicted $m$-shift
is $\Delta m \approx \log_\vp\!\bigl[(m_p/m_e)^2 \cdot (C_e/C_p)\bigr]
\approx -27$, matching the observed shift from $m = -175$ to $m = -202$.

\subsection{Extended dimensionless ratios}
\label{sec:extended}

Beyond the four gauge anchors and gravity, we mapped 43 independent
dimensionless ratios---spanning fermion mass ratios, CKM and PMNS mixing
parameters, anomalous magnetic moments, and energy-scale hierarchies---to
their nearest lattice points.

\begin{table}[H]
\centering
\caption{Combined hit rates across 43 dimensionless ratios.}
\label{tab:combined_hits}
\begin{tabular}{@{}lrrrc@{}}
\toprule
Tolerance & Hits & Rate & Null rate & Enrichment \\
\midrule
$< 1\%$ & 18/43 & 42\% & 15.6\% & $\mathbf{2.68\times}$ \\
$< 3\%$ & 31/43 & 72\% & 39.4\% & $1.83\times$ \\
$< 5\%$ & 38/43 & 88\% & 52.4\% & $1.69\times$ \\
\bottomrule
\end{tabular}
\end{table}

Selected sub-percent hits include:

\begin{table}[H]
\centering
\caption{Tightest lattice hits from the 43-ratio survey (selection).}
\label{tab:tight_hits}
\begin{tabular}{@{}lrrrl@{}}
\toprule
Ratio & Value & $(C,m)$ & $C/\vp^m$ & Rel.~err \\
\midrule
$m_t/m_\tau$        & 97.23   & $(60, -1)$  & 97.08   & $-0.15\%$ \\
$m_Z/\Lambda_{\mathrm{QCD}}$ & 291.69  & $(180, -1)$ & 291.25  & $-0.15\%$ \\
$a_e = (g{-}2)_e/2$ & 0.001160 & $(120, 24)$ & 0.001157 & $-0.20\%$ \\
$|V_{ud}|$          & 0.9737  & $(120, 10)$ & 0.9754  & $+0.20\%$ \\
$m_\tau/m_e$         & 3477    & $(120, -7)$ & 3484    & $+0.20\%$ \\
$\sin\theta_C$      & 0.2253  & $(45, 11)$  & 0.2261  & $+0.37\%$ \\
$m_p/m_e$            & 1836    & $(15, -10)$ & 1845    & $+0.48\%$ \\
$\delta_{\mathrm{CKM}}/\pi$ & 0.364 & $(45, 10)$ & 0.365 & $+0.48\%$ \\
$m_H/m_W$           & 1.558   & $(45, 7)$   & 1.550   & $-0.54\%$ \\
\bottomrule
\end{tabular}
\end{table}

All 43 ratios map to exactly the six $C$~values from the gauge-group
construction: $\{15, 45, 60, 120, 180, 360\}$.  No ratio requires any
value outside this menu.

\subsubsection{CKM and PMNS mixing parameters}

The full CKM and PMNS parameter set (22 quantities) achieves 41\% sub-percent
hits, with notable results including $|V_{ud}| \to (120, 10)$ at $+0.20\%$,
the Cabibbo angle $|V_{us}| \to (45, 11)$ at $+0.81\%$, the CKM CP~phase
$\delta/\pi \to (45, 10)$ at $+0.48\%$, and $\sin^2\theta_{12}^{\mathrm{PMNS}}
\to (60, 11)$ at $-0.82\%$.

\subsection{Lattice-constrained GUT unification}
\label{sec:gut}

Standard grand unification asks where three running couplings converge.  The
lattice adds a stronger constraint: where do all three converge
\emph{onto a single lattice point} $(C, m)$?

Running the three MSSM gauge couplings from $\mZ$ to high energies using
1-loop beta coefficients and searching for the energy $Q$ that minimizes the
maximum deviation of all three $\alpha_i^{-1}(Q)$ from a single lattice value
$C/\vp^m$, we find:

\begin{equation}
\boxed{(C, m) = (15, -1), \quad
\aGUT^{-1} = \frac{15}{\vp^{-1}} = 15\vp \approx 24.271, \quad
\QGUT \approx 1.6 \times 10^{16}~\text{GeV}}
\label{eq:gut}
\end{equation}

with maximum deviation 0.93 across all three couplings.  The RG trajectory
through the lattice shows that $\alpha_2^{-1}$ locks onto $(15, -1)$ first at
$Q \sim 10^{11}$~GeV and remains there for four decades of energy.  At
$Q \sim 5\times 10^{16}$~GeV, two of three couplings simultaneously occupy
$(15, -1)$.

\paragraph{Group-theoretic significance of $C = 15$.}  This value appears
through three independent constructions:
\begin{align}
  15 &= \frac{360}{\dim(\mathrm{SU}(3)) \times h(\mathrm{SU}(3))}
     = \frac{360}{8 \times 3}, \label{eq:c15_su3} \\
  15 &= \frac{360}{\dim(\mathrm{SU}(5))}
     = \frac{360}{24}, \label{eq:c15_su5} \\
  15 &= \dim(\mathrm{SU}(4)), \label{eq:c15_su4}
\end{align}
simultaneously encoding the strong force ($\mathrm{SU}(3)$), the minimal GUT
group ($\mathrm{SU}(5)$), and the Pati--Salam intermediate group
($\mathrm{SU}(4)$).

\paragraph{Significance of $m = -1$.}  The gauge cluster occupies
$m = 0$ to~$+4$; the GUT coupling sits one step into the negative-$m$ side,
at the boundary between the gauge cluster and the gravity realm.  Unification
is the doorstep of gravity on the $m$-axis.

\paragraph{Energy hierarchy in $\vp$-units.}  The scale ratios decompose
cleanly:
\begin{equation}
  \frac{\QGUT}{\mZ} \approx \vp^{70}, \qquad
  \frac{\MPl}{\QGUT} \approx \vp^{12}, \qquad
  \frac{\MPl}{\mZ} \approx \vp^{82},
  \label{eq:hierarchy}
\end{equation}
with the additive relation $70 + 12 = 82$ exact to the nearest integer.  The
Planck-to-GUT gap exponent $n_{\mathrm{gap}} = 12$ equals
$\dim(G_{\SM}) = \dim(U(1)) + \dim(SU(2)) + \dim(SU(3)) = 1 + 3 + 8$.  The
full hierarchy obeys
\begin{equation}
  n_{\mathrm{total}} = \dim(G_{\SM}) \times \vp^4\,.
\end{equation}


% ===================================================================
%  4. STATISTICAL SIGNIFICANCE
% ===================================================================
\section{Statistical Significance}
\label{sec:stats}

\subsection{Null model}
\label{sec:null}

The null hypothesis is that the 43 tested dimensionless ratios bear no
relation to the lattice.  We draw values log-uniformly across 35 decades
(the range spanned by the ratios from $\sim\!10^{-3}$ to $\sim\!10^{38}$)
and compute the probability that each random value falls within a given
tolerance of some lattice point.  The resulting null hit rates are 15.6\%
at $<\!1\%$ tolerance, 39.4\% at $<\!3\%$, and 52.4\% at $<\!5\%$.

\subsection{Binomial significance}
\label{sec:binomial}

Observing 18 hits at $<\!1\%$ out of 43~trials against a null probability of
15.6\% per trial gives a binomial $p$-value:
\begin{equation}
  p = P(X \ge 18 \mid n = 43,\, p_0 = 0.156) \approx 4.5 \times 10^{-5}\,,
\end{equation}
corresponding to roughly 1~in~22{,}000 by chance.

\subsection{Base-uniqueness permutation test}
\label{sec:base_test}

To test whether the base~360 is special, we repeat the 20-ratio independent
prediction test (a subset used for this specific test) with 26 alternative
bases from 60 to~1080.  The top-performing bases by $<\!1\%$ enrichment are
60, 120, 180, 240, 360, 720, and 1080---all divisors or multiples of~360.
Bases outside the 360-family (e.g., 100, 200, 350, 400, 500) consistently
score lower.  The 360-family is collectively special, even though 360 itself
does not maximize the enrichment metric in isolation.

The physical significance of~360 lies not in the enrichment rate but in the
structural properties of its $C$~menu: $C = 15 = 360/\dim(\mathrm{SU}(5))$
enables the GUT unification result (Section~\ref{sec:gut}) and the
group-theoretic connections of Eqs.~\eqref{eq:c15_su3}--\eqref{eq:c15_su4}.

\subsection{Pre-registered out-of-sample predictions}
\label{sec:oos}

Twenty new dimensionless ratios---including neutrino mixing angles, the
Cabibbo angle, the Jarlskog invariant, and the electron anomalous magnetic
moment---were tested after all model rules were frozen.  Standout hits include
$\sin\theta_C \to (45, 11)$ at $+0.37\%$ and
$a_e = (g{-}2)_e/2 \to (120, 24)$ at $-0.20\%$.

\begin{table}[H]
\centering
\caption{In-sample vs.\ out-of-sample hit rates.}
\label{tab:oos}
\begin{tabular}{@{}lrrrr@{}}
\toprule
Tolerance & Original (20) & Pre-registered (20) & Combined (40) & Null \\
\midrule
$< 1\%$ & 30\% & 25\% & 28\% & 15.6\% \\
$< 3\%$ & 75\% & 60\% & 68\% & 39.4\% \\
$< 5\%$ & 85\% & 85\% & 85\% & 52.4\% \\
\bottomrule
\end{tabular}
\end{table}

The enrichment persists out-of-sample at $1.76\times$ ($<\!1\%$).  The slight
decrease from $1.92\times$ (in-sample) to $1.60\times$ (pre-registered) is
the expected regression when moving from in-sample to out-of-sample testing.

\subsection{$C$-value clustering}
\label{sec:clustering}

Monte Carlo simulation (100{,}000 trials, uniform random assignment of
20~items to 6~bins) shows that the raw cluster size---6 ratios in the $C=360$
bin---is not statistically anomalous ($p \approx 0.55$).  The significance of
the $C = 15$ cluster (4~ratios) rests not on the count but on the independent
group-theoretic derivation (Section~\ref{sec:gut}).


% ===================================================================
%  5. THE FINE-TUNING PROBLEMS AS LATTICE INTEGERS
% ===================================================================
\section{The Fine-Tuning Problems as Lattice Integers}
\label{sec:finetuning}

\subsection{The hierarchy problem}
\label{sec:hierarchy}

In the Standard Model, the Higgs mass receives quadratically divergent quantum
corrections from every energy scale up to $\MPl$.  Obtaining
$m_H \sim 125$~GeV requires the bare mass and corrections to cancel to
$\sim\!1$~part in~$10^{34}$.  This is the central motivation for BSM
physics---supersymmetry, extra dimensions, composite Higgs---none of which has
been observed at the LHC.

On the lattice:
\begin{equation}
  \boxed{\frac{\MPl}{m_p} = \vp^{91} \quad (+0.51\%)}
  \label{eq:hierarchy_lattice}
\end{equation}
The 19-order-of-magnitude gulf between gravity and the strong force is 91
$\vp$-steps.  There is no cancellation---only a count.  The question
``why is gravity so weak?''\ becomes ``why~91?''

\subsection{The cosmological constant problem}
\label{sec:cc}

Quantum field theory predicts
$\rho_{\mathrm{vac}} \sim M_{\mathrm{Pl}}^4 \sim 10^{76}$~GeV$^4$.  The
observed value is $\rho_{\mathrm{vac}} \sim 10^{-47}$~GeV$^4$, a discrepancy
of 123 orders of magnitude---the worst prediction in physics.

On the lattice:
\begin{equation}
  \boxed{\frac{\rho_{\mathrm{Pl}}}{\rho_{\mathrm{vac}}} = \vp^{588}
  \quad (+0.01\%)}
  \label{eq:cc_lattice}
\end{equation}
The 123-order-of-magnitude discrepancy is 588 $\vp$-steps at one hundredth
of a percent precision.

\subsection{Number theory of the lattice integers}
\label{sec:number_theory}

Both integers have structured factorizations:
\begin{align}
  91  &= 7 \times 13\,, \\
  588 &= 2^2 \times 3 \times 7^2\,.
\end{align}
They share the factor~7: $\gcd(91, 588) = 7$.  Their ratio is
$588/91 \approx 6.46$, close to $\vp^4 - \vp^0 = 5.85$ and to the Coxeter
number of $\mathrm{SU}(3)$ plus~3.  The Zeckendorf decompositions (unique
sums of non-consecutive Fibonacci numbers) are $91 = 89 + 2$ and
$588 = 377 + 144 + 55 + 8 + 3 + 1$.

The factorization $91 = 7 \times 13$ is notable: both 7 and 13 are primes
adjacent to Fibonacci numbers ($F_6 = 8$, $F_7 = 13$), and 91 is the 13th
triangular number.  The factorization $588 = 12 \times 49$ connects the
cosmological constant to $\dim(G_{\SM}) = 12$ and $7^2 = 49$.

\subsection{Implications}
\label{sec:implications}

Two changes follow from the lattice perspective:
\begin{enumerate}[nosep]
  \item \textbf{Discrete replaces continuous:} in the Standard Model, the
    Higgs mass and vacuum energy are continuous parameters; the fine-tuning
    is that out of uncountably many possible values, nature selects specific
    ones.  On the lattice, the question is: out of the integers near~91
    and~588, why exactly those?  This is a qualitatively smaller question.
  \item \textbf{Common origin:} the two problems, conventionally treated as
    completely independent, are the same \emph{kind} of object on the
    lattice---both are integer $\vp$-powers sharing the factor~7.  This
    suggests a unified mechanism.
\end{enumerate}


% ===================================================================
%  6. FALSIFIABLE PREDICTIONS
% ===================================================================
\section{Falsifiable Predictions}
\label{sec:predictions}

The following predictions are derived from specific lattice address
assignments.  Each is the exact lattice value $C/\vp^m$, compared against the
current experimental measurement and the future experiment that will test it.

\subsection{Tier~1: testable within 5~years}

\paragraph{1. PMNS CP-violation phase.}
\begin{equation}
  \delta_{CP}^{\mathrm{lattice}} = \frac{360}{\vp^{12}} \times 180^\circ
  = 1.118 \times 180^\circ = \mathbf{201.2^\circ}
\end{equation}

\begin{table}[H]
\centering
\begin{tabular}{@{}lrl@{}}
\toprule
 & Value & Source \\
\midrule
Lattice prediction    & $201.2^\circ$ & $(360, 12)$ \\
Current measurement   & $197^\circ \pm 25^\circ$ & T2K/NOvA combined \\
Future precision      & $\pm 5^\circ$ & DUNE ($\sim$2032) \\
\bottomrule
\end{tabular}
\end{table}

\paragraph{2. Atmospheric neutrino mixing.}
\begin{equation}
  \sin^2\theta_{23}^{\mathrm{lattice}} = \frac{180}{\vp^{12}} = \mathbf{0.5590}
\end{equation}

\begin{table}[H]
\centering
\begin{tabular}{@{}lrl@{}}
\toprule
 & Value & Source \\
\midrule
Lattice prediction    & $0.5590$ & $(180, 12)$ \\
Current measurement   & $0.573 \pm 0.020$ & NuFIT~5.2 \\
Future precision      & $\pm 0.01$ & DUNE / Hyper-K ($\sim$2030) \\
\bottomrule
\end{tabular}
\end{table}

\paragraph{3. Reactor neutrino mixing (the most dangerous prediction).}
\begin{equation}
  \sin^2\theta_{13}^{\mathrm{lattice}} = \frac{120}{\vp^{18}} = \mathbf{0.02077}
\end{equation}

\begin{table}[H]
\centering
\begin{tabular}{@{}lrl@{}}
\toprule
 & Value & Source \\
\midrule
Lattice prediction    & $0.02077$ & $(120, 18)$ \\
Current measurement   & $0.02219 \pm 0.00062$ & Daya Bay / RENO \\
Future precision      & $\pm 0.0003$ & JUNO ($\sim$2028) \\
\bottomrule
\end{tabular}
\end{table}

This is the sharpest test: the lattice value is $2.3\sigma$ below the current
measurement (6.4\% discrepancy).  If JUNO confirms $0.0222$, this address
assignment is falsified.  If the value shifts downward, it is a striking
confirmation.

\subsection{Tier~2: testable within 10~years}

\paragraph{4. W~boson mass.}
The lattice assigns $m_H/m_W \to (45, 7)$ at $-0.54\%$, predicting
\begin{equation}
  m_W^{\mathrm{lattice}} = \frac{m_H}{45/\vp^7} = \frac{125.25}{1.5499}
  = \mathbf{80.81~\text{GeV}}\,.
\end{equation}
This sits above both the PDG average ($80.377 \pm 0.012$~GeV, $+0.5\%$)
and the CDF~II anomaly ($80.434 \pm 0.009$~GeV, $+0.5\%$).  FCC-ee will
measure $m_W$ to $\pm 0.5$~MeV ($\sim$2040), resolving this decisively.

\paragraph{5. Proton decay lifetime.}
The lattice-constrained GUT parameters ($\aGUT^{-1} = 15\vp$,
$\QGUT \approx 1.6 \times 10^{16}$~GeV) yield
\begin{equation}
  \tau_p \propto \frac{M_{\mathrm{GUT}}^4}{\alpha_{\mathrm{GUT}}^2\, m_p^5}
  \sim 10^{34\text{--}35}~\text{years}\,,
\end{equation}
squarely in the Hyper-Kamiokande detection window
($\tau_p > 10^{35}$~years for $p \to e^+\pi^0$, operational $\sim$2027).

\paragraph{6. Sum of neutrino masses.}
Cosmological surveys (EUCLID, CMB-S4, DESI) will constrain
$\sum m_\nu$ to $\pm 0.02$~eV.  The lattice predicts normal hierarchy
with $\sum m_\nu \approx 60$~meV, testable once the absolute mass scale is
pinned.

\subsection{Tier~3: structural predictions}

\paragraph{7. New particles must fall on the lattice.}
If any new particle is discovered, its mass ratio to any known particle should
be expressible as $C/\vp^m$ for integer~$m$ and $C \in \{15, 45, 60, 120,
180, 360\}$ with sub-percent accuracy.  This is a standing, open-ended
prediction that every future discovery can test.

\subsection{Tier~4: vacuum spectral structure as an experimental probe}
\label{sec:vacuum}

The lattice assigns a precise address to the cosmological constant:
$\rho_{\mathrm{Pl}}/\rho_{\mathrm{vac}} = \vp^{588}$ at $+0.01\%$
(Section~\ref{sec:cc}).  If the lattice reflects genuine structure, the
vacuum is not featureless---it should exhibit spectral features at frequencies
related by powers of~$\vp$.  This yields three concrete experimental
proposals, all within reach of existing technology:

\paragraph{8. Dynamic Casimir photon production at $\vp$-harmonic
frequencies.}
A SQUID-based parametric oscillator driven at a swept frequency should
produce photon bursts whose rate peaks at frequencies satisfying
$f_n = f_0 \cdot \vp^n$ for integer~$n$, where $f_0$ is anchored to the
dark energy scale $\Lambda_{\mathrm{DE}}^{1/4}/h \approx 560$~GHz.
Off-lattice drive frequencies should yield the standard (featureless)
dynamic Casimir spectrum; $\vp$-harmonic frequencies should yield
anomalously enhanced production.

\paragraph{9. Casimir force at lattice-tuned plate separations.}
The Casimir force between parallel plates depends on the plate separation~$d$.
At separations corresponding to lattice wavelengths
$d_n = c/(f_0\,\vp^n)$, the measured force should deviate from the standard
$1/d^4$ prediction by a calculable amount if the vacuum has $\vp$-periodic
spectral structure.

\paragraph{10. Excess noise in a cryogenic sub-THz cavity.}
A cryogenic cavity with resonant frequency $\sim\!560$~GHz (the dark energy
scale: $\Lambda_{\mathrm{DE}}^{1/4} \approx 2.3$~meV $\approx 560$~GHz)
should detect anomalously high thermal noise compared to a detuned cavity.
The predicted excess is small but nonzero, and the measurement is a direct
probe of the vacuum energy spectral density at the lattice frequency.

\medskip

The vacuum approach is stronger than the gravitational wave conversion route
(inverse Gertsenshtein effect) for two reasons.  First, the vacuum spectral
density at any given frequency, while small, is far larger than the
gravitational wave strain ($h \sim 10^{-25}$); the Casimir effect is
measured routinely.  Second, this route connects directly to the cosmological
constant problem---the single most dramatic lattice result ($\vp^{588}$ at
$0.01\%$).  Detecting $\vp$-periodic structure in vacuum fluctuations would
constitute an experimental probe of the CC problem, a much larger result than
gravitational wave background detection.

These proposals are publishable as an experimental letter because they make
\emph{specific, falsifiable, frequency-targeted predictions} derived from a
framework that already matches 43~dimensionless constants.  A positive result
is a breakthrough; a null result constrains the lattice.  The predictions do
not require the reader to accept the lattice as correct---only that it is
worth testing.

\subsection{Summary}
\label{sec:prediction_summary}

\begin{table}[H]
\centering
\caption{Registered predictions and their experimental tests.}
\label{tab:predictions}
\begin{tabular}{@{}clcccl@{}}
\toprule
\# & Quantity & Lattice & Current & Tension & Experiment \\
\midrule
1 & $\delta_{CP}^{\mathrm{PMNS}}$ & $201.2^\circ$ & $197 \pm 25^\circ$ & $0.2\sigma$ & DUNE \\
2 & $\sin^2\theta_{23}$ & 0.5590 & $0.573 \pm 0.020$ & $0.7\sigma$ & DUNE/HyperK \\
3 & $\sin^2\theta_{13}$ & 0.02077 & $0.02219 \pm 0.00062$ & $\mathbf{2.3\sigma}$ & JUNO \\
4 & $m_W$ & 80.81~GeV & $80.377 \pm 0.012$ & $\mathbf{3.6\sigma}$ & FCC-ee \\
5 & $\tau_p$ & $10^{34\text{--}35}$~yr & $> 10^{34}$~yr & --- & HyperK \\
6 & $\sum m_\nu$ & $\sim$60~meV & $< 120$~meV & --- & EUCLID/S4 \\
7 & new particles & on lattice & --- & --- & LHC/future \\
8 & DCE at $\vp$-harmonics & peaked & featureless & --- & SQUID parametric \\
9 & Casimir at lattice~$d$ & anomalous & $1/d^4$ & --- & precision plates \\
10 & 560~GHz cavity noise & excess & thermal & --- & cryo sub-THz \\
\bottomrule
\end{tabular}
\end{table}

Predictions~3 and~4 are in tension with current data and constitute the
sharpest particle-physics tests.  Predictions~8--10 probe the vacuum
directly and connect to the cosmological constant problem.  Their
outcomes---confirmation or falsification---are scientifically informative
either way.


% ===================================================================
%  7. LIMITATIONS AND OPEN QUESTIONS
% ===================================================================
\section{Limitations and Open Questions}
\label{sec:limitations}

We identify the following limitations, grouped by severity.

\subsection{No dynamical mechanism}
\label{sec:no_mechanism}

The lattice is descriptive: it says $G = C/\vp^m$ but not \emph{why}.  Until
a Lagrangian, variational principle, or constructive derivation produces the
lattice from more fundamental input, the framework remains a phenomenological
observation rather than a theory.  We note five candidate mechanisms as
directions for future investigation, none validated:
\begin{enumerate}[nosep]
  \item \textbf{Discrete scale invariance} at a critical point, with
    preferred scaling factor $\vp$ (motivated by KAM stability and the
    Efimov effect analogy);
  \item \textbf{Conformal fixed point} with Fibonacci fusion rules;
  \item \textbf{Planck-scale quasicrystal} geometry
    (Penrose-tiling-like structure);
  \item \textbf{RG flow with $\vp$-blocking} (coupling constants become
    trapped at $\vp$-spaced fixed points);
  \item \textbf{Maximum entropy on a constrained lattice} (the observed
    configuration maximizes Shannon entropy subject to the lattice constraint).
\end{enumerate}

\subsection{The muon mass outlier}
\label{sec:muon}

The muon mass ratio $m_\mu/m_e = 206.77$ does not fit the lattice at
sub-percent accuracy.  The nearest lattice point, $(120, -1) = 120\vp
\approx 194.2$, is $6.1\%$ off.  This is the most prominent failure among
the charged lepton sector.  Possible resolutions include: a higher-order
lattice effect, a non-minimal $C$~value from representation-dependent
invariants, or a genuine failure mode indicating the lattice structure does
not extend to all mass ratios.  We document this outlier without attempting
to explain it away.

\subsection{Tolerance sensitivity}
\label{sec:tolerance}

At the strict 5\% threshold, all four gauge anchors pass.  At 3\%, the EM
and hypercharge anchors still pass, but the strong ($-1.27\%$) and weak
($-0.73\%$) require no weakening.  The strong coupling passes at 2\%.  If
tolerance is tightened further, the framework's coverage decreases, which is
the expected behavior of a coarse approximation.

\subsection{Gravity ambiguity}
\label{sec:gravity_ambiguity}

The mass-dependence of $\alpha_G(M)$ means that gravity's lattice address is
not unique without specifying the gravitating mass.  This is a feature
(different ``gravity types'' correspond to different mass scales) and a
limitation (without a physical principle selecting the mass anchor, the fit
is under-constrained).

\subsection{Overfitting pathways}
\label{sec:overfitting}

Relaxing the strict $C$~menu or allowing orientation flips (e.g., $\alpha$
instead of $1/\alpha$) rapidly expands the hypothesis space.  The framework's
credibility rests entirely on the frozen contract
(Section~\ref{sec:contract}); any modification requires re-running all
pre-registered tests under the new rules.


% ===================================================================
%  8. CONCLUSION
% ===================================================================
\section{Conclusion}
\label{sec:conclusion}

We have presented a minimal discrete framework in which dimensionless Standard
Model parameters are compared against the lattice $G(C,m) = C/\vp^m$, with
$C$ restricted to six gauge-group-derived values and $m$ forced to be an
integer.  The main findings are:

\begin{enumerate}[nosep]
  \item Across 43 independent dimensionless ratios, the sub-percent hit rate
    is $2.68\times$ the null expectation ($p < 5\times 10^{-5}$), and this
    enrichment persists out-of-sample.
  \item All three MSSM gauge couplings converge on the lattice point
    $(C=15, m=-1)$ with $\aGUT^{-1} = 15\vp \approx 24.27$, where
    $C = 15 = 360/\dim(\mathrm{SU}(5))$.
  \item The Planck-to-proton hierarchy and the cosmological constant
    discrepancy reduce to single lattice integers: $\vp^{91}$ and $\vp^{588}$.
  \item Ten registered predictions are presented.  Particle-physics
    predictions ($\sin^2\theta_{13}$, $\delta_{CP}^{\mathrm{PMNS}}$, $m_W$)
    are testable by DUNE, JUNO, Hyper-Kamiokande, and FCC-ee within the
    coming decade.  Three vacuum spectral predictions---dynamic Casimir
    photon production at $\vp$-harmonic frequencies, anomalous Casimir force
    at lattice-tuned plate separations, and excess noise in a 560~GHz
    cryogenic cavity---provide a direct experimental route to probing the
    cosmological constant problem.
\end{enumerate}

The framework is not a theory: it provides no dynamical mechanism and does not
replace the Standard Model.  It is a phenomenological observation, rigorously
constrained against overfitting, that points toward a discrete organizational
principle in the space of fundamental constants.  The companion open-source
codebase, available at \url{[repository URL]}, reproduces every numerical
result reported in this paper.

Whether the pattern is a deep structural feature of nature or an elaborate
coincidence is a question that experiment will answer.  Predictions~3 and~4
(Table~\ref{tab:predictions}) place the framework at maximal risk in
particle physics within the next five years, while predictions~8--10 offer
a complementary laboratory route through vacuum spectral measurements.


% ===================================================================
%  ACKNOWLEDGMENTS
% ===================================================================
\section*{Acknowledgments}

[To be added.]


% ===================================================================
%  APPENDICES
% ===================================================================
\appendix

\section{Complete lattice address table}
\label{app:address_table}

The full 43-ratio lattice mapping.  Each ratio is independent (not
algebraically derivable from the others within the set).  Ratios are grouped
by sector and sorted by $C$, then~$m$.

\begin{table}[H]
\centering
\footnotesize
\caption{Complete lattice address assignments for all 43 dimensionless ratios.}
\label{tab:full_addresses}
\begin{tabular}{@{}llrrrl@{}}
\toprule
Sector & Ratio & Value & $(C, m)$ & $C/\vp^m$ & Rel.~err \\
\midrule
\multicolumn{6}{@{}l}{\textit{Gauge couplings (5)}} \\
gauge & $1/\alpha_{\mathrm{em}}$ & 137.036 & $(360, 2)$ & 137.508 & $+0.34\%$ \\
gauge & $1/\alpha_s(m_H)$ & 8.867 & $(60, 4)$ & 8.754 & $-1.27\%$ \\
gauge & $1/\alpha_2$ & 28.536 & $(120, 3)$ & 28.328 & $-0.73\%$ \\
gauge & $1/\alpha_1^{\mathrm{GUT}}$ & 59.652 & $(60, 0)$ & 60.000 & $+0.58\%$ \\
gauge & $\alpha/(2\pi)$ & 0.001161 & $(120, 24)$ & 0.001157 & $-0.35\%$ \\
\addlinespace
\multicolumn{6}{@{}l}{\textit{Lepton mass ratios (3)}} \\
lepton & $m_\tau/m_e$ & 3477.2 & $(120, -7)$ & 3484.1 & $+0.20\%$ \\
lepton & $m_\mu/m_e$ & 206.77 & $(120, -1)$ & 194.16 & $-6.10\%$ \\
lepton & $m_\tau/m_\mu$ & 16.82 & $(45, 2)$ & 17.19 & $+2.21\%$ \\
\addlinespace
\multicolumn{6}{@{}l}{\textit{Quark mass ratios (8)}} \\
quark & $m_t/m_u$ & 78{,}527 & $(60, -15)$ & 81{,}840 & $+4.22\%$ \\
quark & $m_t/m_\tau$ & 97.23 & $(60, -1)$ & 97.08 & $-0.15\%$ \\
quark & $m_c/m_s$ & 13.66 & $(60, 3)$ & 14.16 & $+3.72\%$ \\
quark & $m_t/m_b$ & 41.33 & $(180, 3)$ & 42.49 & $+2.81\%$ \\
quark & $m_b/m_\tau$ & 2.352 & $(180, 9)$ & 2.368 & $+0.66\%$ \\
quark & $m_b/m_d$ & 889.4 & $(360, -2)$ & 942.5 & $+5.97\%$ \\
quark & $m_c/m_u$ & 577.3 & $(360, -1)$ & 582.5 & $+0.90\%$ \\
quark & $m_s/m_d$ & 19.79 & $(360, 6)$ & 20.06 & $+1.39\%$ \\
\addlinespace
\multicolumn{6}{@{}l}{\textit{Baryon / meson (4)}} \\
baryon & $m_p/m_e$ & 1836.2 & $(15, -10)$ & 1844.9 & $+0.48\%$ \\
baryon & $(m_n - m_p)/m_e$ & 2.530 & $(45, 6)$ & 2.508 & $-0.89\%$ \\
baryon & $m_p/m_\pi$ & 6.723 & $(120, 6)$ & 6.687 & $-0.52\%$ \\
baryon & $m_n/m_p$ & 1.001 & $(120, 10)$ & 0.976 & $-2.57\%$ \\
\addlinespace
\multicolumn{6}{@{}l}{\textit{Electroweak sector (7)}} \\
EW & $m_t/m_W$ & 2.149 & $(15, 4)$ & 2.188 & $+1.82\%$ \\
EW & $m_H/m_Z$ & 1.374 & $(15, 5)$ & 1.353 & $-1.53\%$ \\
EW & $m_Z/m_t$ & 0.5278 & $(15, 7)$ & 0.5166 & $-2.12\%$ \\
EW & $m_H/m_W$ & 1.558 & $(45, 7)$ & 1.550 & $-0.54\%$ \\
EW & $v_{\mathrm{EW}}/\mZ$ & 2.700 & $(120, 8)$ & 2.554 & $-5.40\%$ \\
EW & $m_Z/m_W$ & 1.134 & $(360, 12)$ & 1.118 & $-1.45\%$ \\
EW & $G_F \mZ^2$ & 0.0970 & $(360, 17)$ & 0.1008 & $+3.94\%$ \\
\addlinespace
\multicolumn{6}{@{}l}{\textit{Scale hierarchies (5)}} \\
hierarchy & $\MPl/\mZ$ & $1.34\times 10^{17}$ & $(120, -72)$ & $1.34\times 10^{17}$ & $-0.10\%$ \\
hierarchy & $\MPl/\QGUT$ & 763.1 & $(180, -3)$ & 762.5 & $-0.07\%$ \\
hierarchy & $\mZ/\Lambda_{\mathrm{QCD}}$ & 291.69 & $(180, -1)$ & 291.25 & $-0.15\%$ \\
hierarchy & $\QGUT/\mZ$ & $1.75\times 10^{14}$ & $(360, -56)$ & $1.82\times 10^{14}$ & $+3.62\%$ \\
hierarchy & $m_p/\Lambda_{\mathrm{QCD}}$ & 3.001 & $(360, 10)$ & 2.927 & $-2.48\%$ \\
\addlinespace
\multicolumn{6}{@{}l}{\textit{CKM mixing (6)}} \\
CKM & $|V_{us}|$ & 0.2243 & $(45, 11)$ & 0.2261 & $+0.81\%$ \\
CKM & $|V_{cb}/V_{ub}|$ & 10.68 & $(45, 3)$ & 10.62 & $-0.54\%$ \\
CKM & $|V_{us}/V_{cb}|$ & 5.498 & $(60, 5)$ & 5.410 & $-1.59\%$ \\
CKM & $|V_{ub}|$ & 0.00382 & $(60, 20)$ & 0.00397 & $+3.83\%$ \\
CKM & $J_{\mathrm{CKM}}$ & $3.18\times 10^{-5}$ & $(60, 30)$ & $3.22\times 10^{-5}$ & $+1.41\%$ \\
CKM & $|V_{cb}|$ & 0.0408 & $(360, 19)$ & 0.0385 & $-5.62\%$ \\
\addlinespace
\multicolumn{6}{@{}l}{\textit{PMNS mixing (3)}} \\
PMNS & $\sin^2\theta_{12}$ & 0.304 & $(60, 11)$ & 0.302 & $-0.82\%$ \\
PMNS & $\sin^2\theta_{23}$ & 0.573 & $(180, 12)$ & 0.559 & $-2.44\%$ \\
PMNS & $\sin^2\theta_{13}$ & 0.02219 & $(120, 18)$ & 0.02077 & $-6.41\%$ \\
\addlinespace
\multicolumn{6}{@{}l}{\textit{Anomalous magnetic moments (2)}} \\
anomalous & $a_e = (g{-}2)_e/2$ & 0.001160 & $(120, 24)$ & 0.001157 & $-0.20\%$ \\
anomalous & $a_\mu = (g{-}2)_\mu/2$ & 0.001166 & $(120, 24)$ & 0.001157 & $-0.73\%$ \\
\bottomrule
\end{tabular}
\end{table}

All 43 ratios map onto exactly the six $C$~values from the gauge-group
construction.  No ratio requires a $C$~value outside the frozen menu.  The
sector counts in parentheses sum to~43.

\section{Reproducibility: CLI commands}
\label{app:cli}

Every numerical result in this paper can be reproduced using the open-source
companion codebase.  After installation (\texttt{pip install -e .}), the
following commands regenerate the key results:

\begin{verbatim}
# Core anchor fits (Table 2)
python -m physics_test.cli strict-gauge-fits --tol 0.05

# RG-within-band validation (Section 3.2)
python -m physics_test.cli oos-predictive-rg --suite v2
python -m physics_test.cli oos-predictive-rg --suite v3

# Gravity fits (Table 3)
python -m physics_test.cli strict-gauge-fits --tol 0.05 \
  --include base/C2_fund,base/T_fund

# GUT unification (Section 3.5)
python -m physics_test.cli gut-compare

# Independent lattice predictions (Section 3.4)
python -m physics_test.cli gut-validate

# Statistical significance (Section 4)
python -m physics_test.cli gut-significance

# Deep structural analysis (full address table)
python -m physics_test.cli gut-deep

# Fine-tuning integers (Section 5)
python -m physics_test.cli gut-deep2

# Extended predictions
python -m physics_test.cli gut-predict2
\end{verbatim}


\section{2-loop RG running coefficients}
\label{app:rg}

The 2-loop coupled running uses the full $3 \times 3$ matrix of 2-loop beta
coefficients~\cite{machacek1983,machacek1984}, integrated numerically via
fourth-order Runge--Kutta:
\begin{equation}
  \frac{d\alpha_i}{d\ln\mu} = \frac{b_i}{2\pi}\alpha_i^2
  + \sum_j \frac{B_{ij}}{8\pi^2}\alpha_i^2\,\alpha_j\,,
\end{equation}
where the 1-loop coefficients are:

\paragraph{Standard Model ($n_f = 6$, one Higgs doublet):}
\begin{equation}
  b_i^{\SM} = \left(\frac{41}{10},\; -\frac{19}{6},\; -7\right)\,.
\end{equation}

\paragraph{MSSM:}
\begin{equation}
  b_i^{\MSSM} = \left(\frac{33}{5},\; 1,\; -3\right)\,.
\end{equation}

The 2-loop coefficient matrix $B_{ij}$ for the SM is:
\begin{equation}
  B^{\SM} = \begin{pmatrix}
    199/50 & 27/10 & 44/5 \\
    9/10   & 35/6  & 12   \\
    11/10  & 9/2   & -26
  \end{pmatrix}\,,
\end{equation}
and for the MSSM:
\begin{equation}
  B^{\MSSM} = \begin{pmatrix}
    199/25 & 27/5  & 88/5 \\
    9/5    & 25    & 24   \\
    11/5   & 9     & 14
  \end{pmatrix}\,.
\end{equation}


% ===================================================================
%  REFERENCES
% ===================================================================
\begin{thebibliography}{99}

\bibitem{georgi1974}
  H.~Georgi and S.~L.~Glashow,
  ``Unity of All Elementary Particle Forces,''
  \emph{Phys.\ Rev.\ Lett.} \textbf{32}, 438 (1974).

\bibitem{pati1974}
  J.~C.~Pati and A.~Salam,
  ``Lepton number as the fourth `color',''
  \emph{Phys.\ Rev.\ D} \textbf{10}, 275 (1974).

\bibitem{bousso2000}
  R.~Bousso and J.~Polchinski,
  ``Quantization of four-form fluxes and dynamical neutralization of the
  cosmological constant,''
  \emph{JHEP} \textbf{0006}, 006 (2000).

\bibitem{barrow1986}
  J.~D.~Barrow and F.~J.~Tipler,
  \emph{The Anthropic Cosmological Principle},
  Oxford University Press (1986).

\bibitem{dirac1937}
  P.~A.~M.~Dirac,
  ``The Cosmological Constants,''
  \emph{Nature} \textbf{139}, 323 (1937).

\bibitem{pdg2024}
  R.~L.~Workman \emph{et al.}\ (Particle Data Group),
  ``Review of Particle Physics,''
  \emph{Prog.\ Theor.\ Exp.\ Phys.} \textbf{2022}, 083C01 (2022),
  and 2024 update.

\bibitem{machacek1983}
  M.~E.~Machacek and M.~T.~Vaughn,
  ``Two-loop renormalization group equations in a general quantum field
  theory.\ I.\ Wave function renormalization,''
  \emph{Nucl.\ Phys.\ B} \textbf{222}, 83 (1983).

\bibitem{machacek1984}
  M.~E.~Machacek and M.~T.~Vaughn,
  ``Two-loop renormalization group equations in a general quantum field
  theory.\ II.\ Yukawa couplings,''
  \emph{Nucl.\ Phys.\ B} \textbf{236}, 221 (1984).

\bibitem{nufit2024}
  I.~Esteban \emph{et al.},
  ``The fate of hints: updated global analysis of three-flavor neutrino
  oscillations,''
  \emph{JHEP} \textbf{09}, 178 (2020); NuFIT~5.2 (2024).

\bibitem{dune2020}
  DUNE Collaboration,
  ``Deep Underground Neutrino Experiment (DUNE), Far Detector Technical
  Design Report, Volume~II,''
  \emph{JINST} \textbf{15}, T08008 (2020).

\bibitem{juno2022}
  JUNO Collaboration,
  ``JUNO physics and detector,''
  \emph{Prog.\ Part.\ Nucl.\ Phys.} \textbf{123}, 103927 (2022).

\bibitem{hyperk2018}
  Hyper-Kamiokande Collaboration,
  ``Hyper-Kamiokande Design Report,''
  arXiv:1805.04163 (2018).

\bibitem{fccee2019}
  FCC Collaboration,
  ``FCC-ee: The Lepton Collider,''
  \emph{Eur.\ Phys.\ J.\ ST} \textbf{228}, 261 (2019).

\bibitem{cohn1964}
  J.~H.~E.~Cohn,
  ``On square Fibonacci numbers,''
  \emph{J.\ London Math.\ Soc.} \textbf{39}, 537 (1964).

\end{thebibliography}

\end{document}
